
\subsection{\S~312.31 Information Amendments}

(a) \textit{Requirement for information amendment.} A sponsor shall report in an information amendment essential information on the IND that is not within the scope of a protocol amendment, IND safety reports, or annual report. Examples of information requiring an information amendment include:

\begin{description}[leftmargin=0.5cm,style=sameline]
\item[(1)] New toxicology, chemistry, or other technical information; or
\item[(2)] A report regarding the discontinuance of a clinical investigation.
\end{description}

(b) \textit{Content and format of an information amendment.} An information amendment is required to bear prominent identification of its contents and to contain: (1) a statement of the nature and purpose of the amendment; (2) an organized submission of the data in a format appropriate for scientific review; and (3) if the sponsor desires FDA to comment, a request for such comment.

(c) \textit{When submitted.} Information amendments to the IND should be submitted as necessary but, to the extent feasible, not more than every 30 days.

\subsubsection{Physical AI Adaptation of \S~312.31}

(d) \textit{Physical AI information amendments.} In addition to the examples in paragraph (a), the following Physical AI-related information requires an information amendment:

\begin{description}[leftmargin=0.8cm,style=sameline]
\item[(i)] New simulation validation data, including results from updated simulation framework versions or new cross-framework validation studies;
\item[(ii)] Updated USL scores for Physical AI systems used in the investigation, whether resulting from new assessment data or changes to the USL scoring methodology;
\item[(iii)] New cybersecurity vulnerability information affecting Physical AI systems used in the investigation, including vendor security advisories and patch deployment status;
\item[(iv)] Changes to the national MCP-PAI standard or the Physical AI system's MCP conformance level that affect the system's interaction with clinical data;
\item[(v)] New published literature or regulatory actions in other jurisdictions related to the safety or performance of Physical AI systems of the same type used in the investigation; and
\item[(vi)] Updates to the digital twin modeling framework, including new calibration data sources, algorithm improvements, or validation results that are not significant enough to require a protocol amendment.
\end{description}

[52 FR 8831, Mar. 19, 1987, as amended at 52 FR 23031, June 17, 1987; 53 FR 1918, Jan. 25, 1988; 67 FR 9585, Mar. 4, 2002. Physical AI adaptation added 17 March 2026.]

\subsection{\S~312.32 IND Safety Reporting}

(a) \textit{Definitions.} The following definitions of terms apply to this section:

\textit{Adverse event} means any untoward medical occurrence associated with the use of a drug in humans, whether or not considered drug related.

\textit{Life-threatening adverse event or life-threatening suspected adverse reaction.} An adverse event or suspected adverse reaction is considered ``life-threatening'' if, in the view of either the investigator or sponsor, its occurrence places the patient or subject at immediate risk of death. It does not include an adverse event or suspected adverse reaction that, had it occurred in a more severe form, might have caused death.

\textit{Serious adverse event or serious suspected adverse reaction.} An adverse event or suspected adverse reaction is considered ``serious'' if, in the view of either the investigator or sponsor, it results in any of the following outcomes: death, a life-threatening adverse event, inpatient hospitalization or prolongation of existing hospitalization, a persistent or significant incapacity or substantial disruption of the ability to conduct normal life functions, or a congenital anomaly/birth defect. Important medical events that may not result in death, be life-threatening, or require hospitalization may be considered serious when, based upon appropriate medical judgment, they may jeopardize the patient or subject and may require medical or surgical intervention to prevent one of the listed outcomes.

\textit{Suspected adverse reaction} means any adverse event for which there is a reasonable possibility that the drug caused the adverse event. For the purposes of IND safety reporting, ``reasonable possibility'' means there is evidence to suggest a causal relationship between the drug and the adverse event.

\textit{Unexpected adverse event or unexpected suspected adverse reaction.} An adverse event or suspected adverse reaction is considered ``unexpected'' if it is not listed in the investigator brochure or is not listed at the specificity or severity that has been observed; or, if an investigator brochure is not required or available, is not consistent with the risk information described in the general investigational plan or elsewhere in the current application, as amended.

(b) \textit{Review of safety information.} The sponsor must promptly review all information relevant to the safety of the drug obtained or otherwise received by the sponsor from foreign or domestic sources, including information derived from any clinical or epidemiological investigations, animal or in vitro studies, reports in the scientific literature, and unpublished scientific papers, as well as reports from foreign regulatory authorities and reports of foreign commercial marketing experience for drugs that are not marketed in the United States.

(c) \textit{IND safety reports.}

\begin{description}[leftmargin=0.5cm,style=sameline]
\item[(1)] The sponsor must notify FDA and all participating investigators in an IND safety report of potential serious risks, from clinical trials or any other source, as soon as possible, but in no case later than 15 calendar days after the sponsor determines that the information qualifies for reporting under paragraph (c)(1)(i), (c)(1)(ii), (c)(1)(iii), or (c)(1)(iv) of this section. In each IND safety report, the sponsor must identify all IND safety reports previously submitted concerning a similar suspected adverse reaction, and must analyze the significance of the suspected adverse reaction in light of previous reports or any other relevant information.

\begin{description}[leftmargin=0.8cm,style=sameline]
\item[(i)] \textit{Serious and unexpected suspected adverse reaction.} The sponsor must report any suspected adverse reaction that is both serious and unexpected, where there is evidence to suggest a causal relationship between the drug and the adverse event.
\item[(ii)] \textit{Findings from other studies.} The sponsor must report any findings from epidemiological studies, pooled analysis, or clinical studies that suggest a significant risk in humans exposed to the drug.
\item[(iii)] \textit{Findings from animal or in vitro testing.} The sponsor must report any findings from animal or in vitro testing that suggest a significant risk in humans exposed to the drug, such as reports of mutagenicity, teratogenicity, or carcinogenicity.
\item[(iv)] \textit{Increased rate of occurrence of serious suspected adverse reactions.} The sponsor must report any clinically important increase in the rate of a serious suspected adverse reaction over that listed in the protocol or investigator brochure.
\item[(v)] \textit{Submission of IND safety reports.} The sponsor must submit each IND safety report in a narrative format or on FDA Form 3500A or in an electronic format that FDA can process, review, and archive.
\end{description}

\item[(2)] \textit{Unexpected fatal or life-threatening suspected adverse reaction reports.} The sponsor must also notify FDA of any unexpected fatal or life-threatening suspected adverse reaction as soon as possible but in no case later than 7 calendar days after the sponsor's initial receipt of the information.

\item[(3)] \textit{Reporting format or frequency.} FDA may require a sponsor to submit IND safety reports in a format or at a frequency different than that required under this paragraph.

\item[(4)] \textit{Investigations of marketed drugs.} A sponsor of a clinical study of a drug marketed or approved in the United States that is conducted under an IND is required to submit IND safety reports for suspected adverse reactions observed in the clinical study.

\item[(5)] \textit{Reporting study endpoints.} Study endpoints (e.g., mortality or major morbidity) must be reported to FDA as described in the protocol and ordinarily would not be reported under paragraph (c) of this section, unless a serious and unexpected adverse event occurs for which there is evidence suggesting a causal relationship between the drug and the event.
\end{description}

(d) \textit{Followup.}

\begin{description}[leftmargin=0.5cm,style=sameline]
\item[(1)] The sponsor must promptly investigate all safety information it receives.
\item[(2)] Relevant followup information to an IND safety report must be submitted as soon as the information is available and must be identified as a ``Followup IND Safety Report.''
\item[(3)] If the results of a sponsor's investigation show that an adverse event not initially determined to be reportable under paragraph (c) of this section is so reportable, the sponsor must report it as soon as possible, but in no case later than 15 calendar days after the determination is made.
\end{description}

(e) \textit{Disclaimer.} A safety report or other information submitted by a sponsor under this part does not necessarily reflect a conclusion by the sponsor or FDA that the report or information constitutes an admission that the drug caused or contributed to an adverse event.

\subsubsection{Physical AI Adaptation of \S~312.32}

(f) \textit{Physical AI safety event definitions.} In addition to the definitions in paragraph (a), the following definitions apply to safety reporting for investigations involving Physical AI systems:

\textit{Physical AI adverse event} means any untoward occurrence associated with the operation of a Physical AI system during a clinical investigation involving an investigational drug, whether or not considered to be caused by the Physical AI system. Physical AI adverse events include robotic malfunction, AI prediction error, sensor failure, communication loss, unauthorized system access, digital twin divergence, and any event requiring activation of the emergency stop.

\textit{Serious Physical AI adverse event} means a Physical AI adverse event that results in, or has the potential to result in, any of the outcomes listed in the definition of serious adverse event in paragraph (a), or that results in unplanned interruption of an investigational drug administration procedure requiring clinical intervention, uncontrolled release of the investigational drug, or compromise of the Physical AI system's sterile field during a surgical procedure.

(g) \textit{Physical AI safety reporting requirements.} The sponsor must include Physical AI adverse events in IND safety reports as follows:

\begin{description}[leftmargin=0.8cm,style=sameline]
\item[(i)] \textit{Serious Physical AI adverse events.} Any serious Physical AI adverse event must be reported under the same timelines and procedures as paragraph (c)(1) of this section. The report must include the Physical AI system identification, the nature of the malfunction or error, the system telemetry data for the period surrounding the event, the emergency stop activation status, and the patient outcome.
\item[(ii)] \textit{Physical AI system-drug interaction events.} Any event in which a Physical AI system error directly or indirectly results in incorrect drug administration (wrong dose, wrong route, wrong timing, wrong patient) must be reported as an IND safety report regardless of whether the event meets the criteria for a serious adverse event.
\item[(iii)] \textit{Cybersecurity incidents.} Any cybersecurity incident affecting a Physical AI system during a clinical investigation, including unauthorized access attempts, data integrity violations, and communication interceptions, must be reported in an IND safety report within 15 calendar days of detection. If the incident results in or has the potential to result in patient harm, the 7-calendar-day reporting timeline for unexpected fatal or life-threatening events applies.
\item[(iv)] \textit{AI model performance degradation.} The sponsor must report any sustained degradation in AI/ML model performance below the pre-specified acceptance criteria documented in the PCCP or the Physical AI System Description. Sustained degradation is defined as performance below acceptance criteria for three or more consecutive procedures or for a cumulative period of 24 hours, whichever occurs first.
\item[(v)] \textit{Sim-to-real divergence.} The sponsor must report any instance in which the observed behavior of a Physical AI system during a clinical procedure diverges from the simulated behavior by more than twice the validated sim-to-real gap tolerance (i.e., trajectory deviation greater than 4mm or force discrepancy greater than 1.0N).
\item[(vi)] \textit{Digital twin discrepancy.} The sponsor must report any instance in which a digital twin model used for treatment planning or intraoperative guidance produces predictions that diverge from observed clinical outcomes by more than the pre-specified accuracy thresholds documented in the protocol.
\end{description}

(h) \textit{Physical AI safety review.} In addition to the review requirements of paragraph (b), the sponsor must establish a Physical AI Safety Review Committee (or equivalent function) responsible for the ongoing review and evaluation of Physical AI adverse events. The committee shall include members with expertise in robotic systems engineering, AI/ML, clinical oncology, and cybersecurity. The committee shall meet at intervals not exceeding 90 days during active clinical investigations to review Physical AI safety data and to recommend modifications to Physical AI system operation, human oversight requirements, or USL threshold criteria.

(i) \textit{Hash-chained audit trail preservation.} For all Physical AI adverse events reported under this section, the sponsor must preserve the complete hash-chained audit trail for the Physical AI system covering the period from 24 hours before the event to 72 hours after the event. The audit trail must be maintained in accordance with the electronic records requirements of 21 CFR Part 11 and must be made available to FDA upon request.

[75 FR 59961, Sept. 29, 2010. Physical AI adaptation added 17 March 2026.]

\subsection{\S~312.33 Annual Reports}

A sponsor shall within 60 days of the anniversary date that the IND went into effect, submit a brief report of the progress of the investigation that includes:

(a) \textit{Individual study information.} A brief summary of the status of each study in progress and each study completed during the previous year, including: (1) the title of the study, its purpose, a brief statement identifying the patient population, and whether the study is completed; (2) the total number of subjects initially planned, the number entered to date tabulated by age group, gender, and race, the number whose participation was completed as planned, and the number who dropped out for any reason; and (3) if the study has been completed or interim results are known, a brief description of any available study results.

(b) \textit{Summary information.} Information obtained during the previous year's clinical and nonclinical investigations, including: (1) a narrative or tabular summary of the most frequent and most serious adverse experiences by body system; (2) a summary of all IND safety reports submitted during the past year; (3) a list of subjects who died during participation in the investigation, with the cause of death for each; (4) a list of subjects who dropped out in association with any adverse experience; (5) a brief description of pertinent findings regarding the drug's actions; (6) a list of preclinical studies completed or in progress during the past year; and (7) a summary of any significant manufacturing or microbiological changes made during the past year.

(c) A description of the general investigational plan for the coming year to replace that submitted 1 year earlier.

(d) If the investigator brochure has been revised, a description of the revision and a copy of the new brochure.

(e) A description of any significant Phase 1 protocol modifications made during the previous year and not previously reported to the IND in a protocol amendment.

(f) A brief summary of significant foreign marketing developments with the drug during the past year.

(g) If desired by the sponsor, a log of any outstanding business with respect to the IND for which the sponsor requests or expects a reply, comment, or meeting.

\subsubsection{Physical AI Adaptation of \S~312.33}

(h) \textit{Physical AI system performance summary.} In addition to the information required by paragraphs (a) through (g), the annual report for investigations involving Physical AI systems shall include:

\begin{description}[leftmargin=0.8cm,style=sameline]
\item[(i)] A summary of all Physical AI adverse events reported during the previous year, organized by system type, event category (malfunction, AI error, cybersecurity, digital twin discrepancy), and patient outcome;
\item[(ii)] Updated USL scores for each Physical AI system used in the investigation, with a comparison to the scores reported in the initial IND submission or the most recent annual report, and an explanation of any changes;
\item[(iii)] A summary of AI/ML model updates implemented during the year, distinguishing between updates within PCCP scope and updates requiring protocol amendments, and including performance metrics before and after each update;
\item[(iv)] Updated simulation validation results, including any new cross-framework validation studies, updated sim-to-real gap measurements, and the results of any re-validation prompted by Physical AI system modifications;
\item[(v)] A summary of digital twin model performance, including calibration accuracy, prediction reliability, and any instances of clinically significant divergence from observed outcomes;
\item[(vi)] A summary of cybersecurity events affecting Physical AI systems, including vulnerability disclosures, patch deployments, and incident investigations;
\item[(vii)] A summary of the Physical AI Safety Review Committee's meetings, findings, and recommendations during the year; and
\item[(viii)] An updated Physical AI investigational plan for the coming year, including planned system upgrades, new Physical AI system introductions, and changes to autonomy levels or human oversight requirements.
\end{description}

[52 FR 8831, Mar. 19, 1987, as amended at 52 FR 23031, June 17, 1987; 63 FR 6862, Feb. 11, 1998; 67 FR 9585, Mar. 4, 2002. Physical AI adaptation added 17 March 2026.]

\subsection{\S~312.38 Withdrawal of an IND}

(a) At any time a sponsor may withdraw an effective IND without prejudice.

(b) If an IND is withdrawn, FDA shall be so notified, all clinical investigations conducted under the IND shall be ended, all current investigators notified, and all stocks of the drug returned to the sponsor or otherwise disposed of at the request of the sponsor in accordance with \S~312.59.

(c) If an IND is withdrawn because of a safety reason, the sponsor shall promptly so inform FDA, all participating investigators, and all reviewing Institutional Review Boards, together with the reasons for such withdrawal.

\subsubsection{Physical AI Adaptation of \S~312.38}

(d) \textit{Physical AI system decommissioning upon IND withdrawal.} When an IND is withdrawn and the investigation involved Physical AI systems, the sponsor shall, in addition to the requirements of paragraphs (a) through (c):

\begin{description}[leftmargin=0.8cm,style=sameline]
\item[(i)] Ensure that all Physical AI systems are safely deactivated and removed from clinical service in accordance with the decommissioning procedures specified in the Physical AI System Description;
\item[(ii)] Preserve all Physical AI system logs, telemetry data, hash-chained audit trails, simulation validation records, and digital twin models for the retention period specified in \S~312.57;
\item[(iii)] Securely erase all patient-specific data from Physical AI system local storage in compliance with HIPAA Safe Harbor de-identification requirements, while maintaining de-identified copies in the sponsor's central records;
\item[(iv)] Notify all sites of the Physical AI system decommissioning requirements and verify completion of decommissioning activities at each site within 30 days of the IND withdrawal; and
\item[(v)] If the withdrawal is for a safety reason related to the Physical AI system, provide FDA with a detailed technical analysis of the Physical AI system failure modes and the contributing factors, within 60 days of the withdrawal.
\end{description}

[52 FR 8831, Mar. 19, 1987, as amended at 52 FR 23031, June 17, 1987; 67 FR 9586, Mar. 4, 2002. Physical AI adaptation added 17 March 2026.]

\clearpage

\section{SUBPART C --- ADMINISTRATIVE ACTIONS}

\subsection{\S~312.40 General Requirements for Use of an Investigational New Drug in a Clinical Investigation}

(a) An investigational new drug may be used in a clinical investigation if the following conditions are met:

\begin{description}[leftmargin=0.5cm,style=sameline]
\item[(1)] The sponsor of the investigation submits an IND for the drug to FDA; the IND is in effect under paragraph (b) of this section; and the sponsor complies with all applicable requirements in this part and parts 50 and 56 with respect to the conduct of the clinical investigations; and
\item[(2)] Each participating investigator conducts his or her investigation in compliance with the requirements of this part and parts 50 and 56.
\end{description}

(b) An IND goes into effect:

\begin{description}[leftmargin=0.5cm,style=sameline]
\item[(1)] Thirty days after FDA receives the IND, unless FDA notifies the sponsor that the investigations described in the IND are subject to a clinical hold under \S~312.42; or
\item[(2)] On earlier notification by FDA that the clinical investigations in the IND may begin. FDA will notify the sponsor in writing of the date it receives the IND.
\end{description}

(c) A sponsor may ship an investigational new drug to investigators named in the IND:

\begin{description}[leftmargin=0.5cm,style=sameline]
\item[(1)] Thirty days after FDA receives the IND; or
\item[(2)] On earlier FDA authorization to ship the drug.
\end{description}

(d) An investigator may not administer an investigational new drug to human subjects until the IND goes into effect under paragraph (b) of this section.

\subsubsection{Physical AI Adaptation of \S~312.40}

(e) \textit{Physical AI system readiness requirements.} In addition to the conditions in paragraph (a), when a clinical investigation involves Physical AI systems, the following conditions must also be met before the investigational drug may be administered through or monitored by a Physical AI system:

\begin{description}[leftmargin=0.8cm,style=sameline]
\item[(i)] The Physical AI System Description required under \S~312.23(g) has been submitted as part of the IND and has not been the subject of a clinical hold under \S~312.42;
\item[(ii)] Each Physical AI system to be used in the investigation has met or exceeded the minimum USL threshold applicable to its procedure type, as specified in \S~312.1(e), or has been granted a USL threshold waiver under \S~312.10(d);
\item[(iii)] The pre-procedure safety matrix for each Physical AI system has been verified at the investigational site, including valid authorization tokens, robot readiness verification, simulation validation confirmation, and environmental checks;
\item[(iv)] All investigators and subinvestigators who will operate or supervise the Physical AI system have completed the training requirements specified in the Physical AI System Description and have demonstrated competency through assessments documented in the investigator's file; and
\item[(v)] The site's MCP server infrastructure, if applicable, has been configured and tested to the conformance level required by the protocol.
\end{description}

(f) \textit{Physical AI system deployment.} A sponsor may deploy Physical AI system components (hardware, software, simulation environments) to investigational sites named in the IND under the same timelines as paragraph (c) of this section for drug shipment. Physical AI system components shall not be used in clinical procedures involving investigational drugs until the IND is in effect and all conditions of paragraph (e) have been verified at the site.

[52 FR 8831, Mar. 19, 1987. Physical AI adaptation added 17 March 2026.]

\subsection{\S~312.41 Comment and Advice on an IND}

(a) FDA may at any time during the course of the investigation communicate with the sponsor orally or in writing about deficiencies in the IND or about FDA's need for more data or information.

(b) On the sponsor's request, FDA will provide advice on specific matters relating to an IND. Examples of such advice may include advice on the adequacy of technical data to support an investigational plan, on the design of a clinical trial, and on whether proposed investigations are likely to produce the data and information that is needed to meet requirements for a marketing application.

(c) Unless the communication is accompanied by a clinical hold order under \S~312.42, FDA communications with a sponsor under this section are solely advisory and do not require any modification in the planned or ongoing clinical investigations or response to the agency.

\subsubsection{Physical AI Adaptation of \S~312.41}

(d) \textit{Physical AI advisory communications.} FDA may communicate with the sponsor regarding deficiencies in the Physical AI System Description, simulation validation data, cybersecurity assessment, or other Physical AI-specific components of the IND. The sponsor may request FDA advice on the adequacy of Physical AI system documentation, the appropriateness of USL thresholds for specific clinical procedures, the acceptability of simulation validation methodologies, and the design of human oversight protocols for Physical AI-assisted drug administration.

(e) \textit{Physical AI system demonstrations.} At FDA's request or the sponsor's initiative, the sponsor may arrange for a demonstration of the Physical AI system's capabilities in a non-clinical setting as part of the advisory process. Such demonstrations do not constitute clinical investigations and are not subject to the IND requirements of this part. Demonstration protocols, results, and any FDA feedback shall be documented and maintained in the sponsor's IND file.

[52 FR 8831, Mar. 19, 1987, as amended at 52 FR 23031, June 17, 1987; 67 FR 9586, Mar. 4, 2002. Physical AI adaptation added 17 March 2026.]

\subsection{\S~312.42 Clinical Holds and Requests for Modification}

(a) \textit{General.} A clinical hold is an order issued by FDA to the sponsor to delay a proposed clinical investigation or to suspend an ongoing investigation. The clinical hold order may apply to one or more of the investigations covered by an IND. When a proposed study is placed on clinical hold, subjects may not be given the investigational drug. When an ongoing study is placed on clinical hold, no new subjects may be recruited to the study and placed on the investigational drug; patients already in the study should be taken off therapy involving the investigational drug unless specifically permitted by FDA in the interest of patient safety.

(b) \textit{Grounds for imposition of clinical hold.}

\begin{description}[leftmargin=0.5cm,style=sameline]
\item[(1)] \textit{Clinical hold of a Phase 1 study under an IND.} FDA may place a proposed or ongoing Phase 1 investigation on clinical hold if it finds that:
\begin{description}[leftmargin=0.8cm,style=sameline]
\item[(i)] Human subjects are or would be exposed to an unreasonable and significant risk of illness or injury;
\item[(ii)] The clinical investigators named in the IND are not qualified by reason of their scientific training and experience to conduct the investigation described in the IND;
\item[(iii)] The investigator brochure is misleading, erroneous, or materially incomplete; or
\item[(iv)] The IND does not contain sufficient information required under \S~312.23 to assess the risks to subjects of the proposed studies.
\item[(v)] The IND is for the study of an investigational drug intended to treat a life-threatening disease or condition that affects both genders, and men or women with reproductive potential who have the disease or condition being studied are excluded from eligibility because of a risk or potential risk from use of the investigational drug of reproductive toxicity or developmental toxicity, except under special circumstances as described in the regulation.
\end{description}

\item[(2)] \textit{Clinical hold of a Phase 2 or 3 study under an IND.} FDA may place a proposed or ongoing Phase 2 or 3 investigation on clinical hold if it finds that: (i) any of the conditions in paragraphs (b)(1)(i) through (b)(1)(v) of this section apply; or (ii) the plan or protocol for the investigation is clearly deficient in design to meet its stated objectives.

\item[(3)] \textit{Clinical hold of an expanded access IND or expanded access protocol.} FDA may place an expanded access IND or expanded access protocol on clinical hold if the pertinent criteria in subpart I of this part for permitting the expanded access use to begin are not satisfied, the expanded access submission does not comply with the requirements for expanded access submissions in subpart I, or the pertinent criteria for permitting the expanded access are no longer satisfied.

\item[(4)] \textit{Clinical hold of any study that is not designed to be adequate and well-controlled.} FDA may place a proposed or ongoing investigation that is not designed to be adequate and well-controlled on clinical hold if it finds that any of the conditions in paragraph (b)(1) or (b)(2) apply, or if there is reasonable evidence that the investigation is impeding enrollment in or interfering with an adequate and well-controlled study, or if insufficient quantities of the drug exist, or if the drug has been studied in adequate and well-controlled investigations that strongly suggest lack of effectiveness, among other grounds specified in the regulation.

\item[(5)] \textit{Clinical hold of any investigation involving an exception from informed consent under \S~50.24 of this chapter.} FDA may place such an investigation on clinical hold if any of the conditions in paragraphs (b)(1) or (b)(2) apply, or if the pertinent criteria in \S~50.24 are not submitted or not satisfied.

\item[(6)] \textit{Clinical hold of any investigation involving an exception from informed consent under \S~50.23(d) of this chapter.} FDA may place such an investigation on clinical hold if any of the conditions in paragraphs (b)(1) or (b)(2) apply, or if a determination by the President to waive the prior consent requirement has not been made.
\end{description}

(c) \textit{Discussion of deficiency.} Whenever FDA concludes that a deficiency exists in a clinical investigation that may be grounds for the imposition of clinical hold, FDA will, unless patients are exposed to immediate and serious risk, attempt to discuss and satisfactorily resolve the matter with the sponsor before issuing the clinical hold order.

(d) \textit{Imposition of clinical hold.} The clinical hold order may be made by telephone or other means of rapid communication or in writing. The clinical hold order will identify the studies under the IND to which the hold applies, and will briefly explain the basis for the action. As soon as possible, and no more than 30 days after imposition, the Division Director will provide the sponsor a written explanation of the basis for the hold.

(e) \textit{Resumption of clinical investigations.} An investigation may only resume after FDA has notified the sponsor that the investigation may proceed. Resumption will be authorized when the sponsor corrects the deficiency(ies) previously cited or otherwise satisfies the agency that the investigation(s) can proceed. If a sponsor requests in writing that the clinical hold be removed and submits a complete response, FDA shall respond in writing within 30 calendar days.

(f) \textit{Appeal.} If the sponsor disagrees with the reasons cited for the clinical hold, the sponsor may request reconsideration in accordance with \S~312.48.

(g) \textit{Conversion of IND on clinical hold to inactive status.} If all investigations covered by an IND remain on clinical hold for 1 year or more, the IND may be placed on inactive status by FDA under \S~312.45.

\subsubsection{Physical AI Adaptation of \S~312.42}

(h) \textit{Physical AI grounds for clinical hold.} In addition to the grounds specified in paragraphs (b)(1) through (b)(6), FDA may place a proposed or ongoing clinical investigation on clinical hold based on Physical AI system deficiencies, including:

\begin{description}[leftmargin=0.8cm,style=sameline]
\item[(i)] \textit{Robotic safety failure.} The Physical AI system has demonstrated a pattern of safety-critical failures, including unplanned emergency stop activations exceeding the rate specified in the protocol, collision events resulting in or potentially resulting in patient injury, or force limit exceedances during clinical procedures;
\item[(ii)] \textit{AI model degradation.} One or more AI/ML models embedded in the Physical AI system have demonstrated sustained performance degradation below the pre-specified acceptance criteria, and the sponsor has not implemented effective corrective action within the timeframe specified in the PCCP or Physical AI System Description;
\item[(iii)] \textit{Simulation-reality divergence.} The observed behavior of the Physical AI system during clinical procedures diverges from the validated simulation predictions by more than twice the established sim-to-real gap tolerance, indicating that the pre-clinical simulation validation no longer adequately predicts the system's clinical behavior;
\item[(iv)] \textit{Cybersecurity compromise.} The Physical AI system has been the subject of a confirmed cybersecurity breach that compromises or has the potential to compromise patient safety, data integrity, or the security of the investigational drug administration pathway;
\item[(v)] \textit{USL score decline.} The Physical AI system's USL score has declined below the minimum threshold required for its procedure type, whether due to identified deficiencies in simulation switching, AI integration, cross-robot sharing, or clinical trial collaboration capabilities;
\item[(vi)] \textit{Inadequate Physical AI System Description.} The Physical AI System Description submitted under \S~312.23(g) does not contain sufficient information to assess the risks to subjects posed by the Physical AI system, including inadequate simulation validation data, incomplete cybersecurity assessment, or insufficient human oversight protocols;
\item[(vii)] \textit{Digital twin failure.} A digital twin model used for treatment planning or intraoperative guidance has produced clinically significant erroneous predictions that resulted in or could have resulted in inappropriate drug administration, and the sponsor has not demonstrated adequate corrective action; and
\item[(viii)] \textit{Human oversight failure.} The investigation demonstrates a pattern of inadequate human oversight of Physical AI systems, including instances where operators failed to intervene when system performance indicators warranted intervention, or where the operator-to-system ratio was insufficient to maintain effective supervisory control.
\end{description}

(i) \textit{Physical AI clinical hold procedures.} When a clinical hold is imposed on the basis of Physical AI system deficiencies under paragraph (h):

\begin{description}[leftmargin=0.8cm,style=sameline]
\item[(i)] The sponsor shall immediately cease all Physical AI system operations in clinical procedures involving the investigational drug at all affected sites. Patients currently undergoing a Physical AI-assisted procedure at the time of the hold order shall have the procedure completed using manual methods under the supervising investigator's clinical judgment, unless FDA specifically authorizes completion of the procedure using the Physical AI system in the interest of patient safety;
\item[(ii)] The sponsor shall place all affected Physical AI systems in a safe standby mode and preserve all system logs, telemetry data, and hash-chained audit trails from the period preceding the clinical hold;
\item[(iii)] The clinical hold order shall specify whether the hold applies to all Physical AI systems under the IND or only to specific system types, models, or sites; and
\item[(iv)] To resume Physical AI system operations, the sponsor must demonstrate that the Physical AI deficiency has been corrected and that the system has been re-validated through the simulation validation process described in \S~312.21(d), unless FDA determines that a lesser level of re-validation is sufficient based on the nature of the deficiency.
\end{description}

[52 FR 8831, Mar. 19, 1987, as amended at 52 FR 19477, May 22, 1987; 57 FR 13249, Apr. 15, 1992; 61 FR 51530, Oct. 2, 1996; 63 FR 68678, Dec. 14, 1998; 64 FR 54189, Oct. 5, 1999; 65 FR 34971, June 1, 2000; 74 FR 40942, Aug. 13, 2009. Physical AI adaptation added 17 March 2026.]

\subsection{\S~312.44 Termination}

(a) \textit{General.} This section describes the procedures under which FDA may terminate an IND. If an IND is terminated, the sponsor shall end all clinical investigations conducted under the IND and recall or otherwise provide for the disposition of all unused supplies of the drug. A termination action may be based on deficiencies in the IND or in the conduct of an investigation under an IND. Except as provided in paragraph (d) of this section, a termination shall be preceded by a proposal to terminate by FDA and an opportunity for the sponsor to respond. FDA will, in general, only initiate an action under this section after first attempting to resolve differences informally or through the clinical hold procedures described in \S~312.42.

(b) \textit{Grounds for termination.}

\begin{description}[leftmargin=0.5cm,style=sameline]
\item[(1)] \textit{Phase 1.} FDA may propose to terminate an IND during Phase 1 if it finds that:
\begin{description}[leftmargin=0.8cm,style=sameline]
\item[(i)] Human subjects would be exposed to an unreasonable and significant risk of illness or injury.
\item[(ii)] The IND does not contain sufficient information required under \S~312.23 to assess the safety to subjects.
\item[(iii)] The methods, facilities, and controls used for manufacturing, processing, and packing of the investigational drug are inadequate.
\item[(iv)] The clinical investigations are being conducted in a manner substantially different than that described in the protocols submitted in the IND.
\item[(v)] The drug is being promoted or distributed for commercial purposes not justified by the requirements of the investigation.
\item[(vi)] The IND contains an untrue statement of a material fact or omits material information required by this part.
\item[(vii)] The sponsor fails promptly to investigate and inform FDA and all investigators of serious and unexpected adverse experiences in accordance with \S~312.32 or fails to make any other report required under this part.
\item[(viii)] The sponsor fails to submit an accurate annual report in accordance with \S~312.33.
\item[(ix)] The sponsor fails to comply with any other applicable requirement of this part, part 50, or part 56.
\item[(x)] The IND has remained on inactive status for 5 years or more.
\item[(xi)] The sponsor fails to delay a proposed investigation or to suspend an ongoing investigation that has been placed on clinical hold under \S~312.42(b)(4).
\end{description}

\item[(2)] \textit{Phase 2 or 3.} FDA may propose to terminate an IND during Phase 2 or Phase 3 if FDA finds that: (i) any of the conditions in paragraphs (b)(1)(i) through (b)(1)(xi) apply; or (ii) the investigational plan or protocol(s) is not reasonable as a bona fide scientific plan to determine whether or not the drug is safe and effective for use; or (iii) there is convincing evidence that the drug is not effective for the purpose for which it is being investigated.

\item[(3)] FDA may propose to terminate a treatment IND if it finds that: (i) any of the conditions in paragraphs (b)(1)(i) through (x) apply; or (ii) any of the conditions in \S~312.42(b)(3) apply.
\end{description}

(c) \textit{Opportunity for sponsor response.} If FDA proposes to terminate an IND, FDA will notify the sponsor in writing and invite correction or explanation within 30 days. If the sponsor does not respond, the IND shall be terminated. If the sponsor responds but FDA does not accept the explanation, the sponsor shall be informed and provided an opportunity for a regulatory hearing under part 16.

(d) \textit{Immediate termination of IND.} Notwithstanding paragraphs (a) through (c), if at any time FDA concludes that continuation of the investigation presents an immediate and substantial danger to the health of individuals, the agency shall immediately terminate the IND by written notice to the sponsor from the Director of the Center for Drug Evaluation and Research or the Director of the Center for Biologics Evaluation and Research. An IND so terminated is subject to reinstatement by the Director on the basis of additional submissions that eliminate such danger.

\subsubsection{Physical AI Adaptation of \S~312.44}

(e) \textit{Physical AI grounds for termination.} In addition to the grounds specified in paragraph (b), FDA may propose to terminate an IND based on Physical AI system deficiencies that present an ongoing and unresolvable risk to human subjects, including:

\begin{description}[leftmargin=0.8cm,style=sameline]
\item[(i)] Repeated Physical AI safety failures that have not been corrected despite clinical hold and sponsor response, demonstrating a fundamental inadequacy in the Physical AI system's design, manufacture, or operation for the intended clinical use;
\item[(ii)] A confirmed cybersecurity compromise of the Physical AI system that results in unauthorized modification of drug administration parameters, patient data exposure, or loss of system integrity, and for which the sponsor cannot demonstrate adequate remediation;
\item[(iii)] Systematic failure of the sponsor to report Physical AI adverse events in accordance with \S~312.32(g), to maintain Physical AI system logs and audit trails in accordance with \S~312.57, or to comply with the Physical AI-specific requirements of this part; and
\item[(iv)] Evidence that the Physical AI system's simulation validation was materially deficient, misleading, or falsified, such that the pre-clinical safety basis for the investigation is fundamentally compromised.
\end{description}

(f) \textit{Immediate termination for Physical AI emergencies.} FDA may immediately terminate an IND under paragraph (d) when a Physical AI system malfunction presents an immediate and substantial danger to patient health. Such circumstances include a Physical AI system that has caused or is likely to cause serious patient injury due to uncontrolled robotic movement, incorrect drug delivery at dangerous doses, or loss of emergency stop functionality. Upon immediate termination, the sponsor shall decommission all Physical AI systems in accordance with \S~312.38(d).

[52 FR 8831, Mar. 19, 1987, as amended at 52 FR 23031, June 17, 1987; 55 FR 11579, Mar. 29, 1990; 57 FR 13249, Apr. 15, 1992; 67 FR 9586, Mar. 4, 2002. Physical AI adaptation added 17 March 2026.]
